\section*{Введение}
\addcontentsline{toc}{section}{Введение}

Планирование групповых поездок в повседневной практике чаще всего выполняется в нескольких цифровых каналах одновременно: обсуждения ведутся в мессенджерах, маршрут фиксируется в картографических сервисах, расходы учитываются в отдельных приложениях, а погодные данные и справочная информация ищутся во внешних источниках. Такая фрагментация приводит к дублированию действий, потере контекста принятых решений и повышенной координационной нагрузке на организатора поездки [8].

Актуальность темы обусловлена потребностью в едином мобильном инструменте, который поддерживает полный цикл совместного планирования: от формулирования идей и обсуждений до формирования маршрута по дням и прозрачного учета расходов. Для малых групп пользователей критична не только доступность отдельных функций, но и их согласованная работа в общем сценарии [7][8].

Объектом исследования является процесс совместного планирования поездок в малых группах.

Предметом исследования являются пользовательские сценарии, функциональные требования и проектные решения мобильного приложения, обеспечивающего совместную работу участников поездки в едином информационном пространстве.

Цель выпускной квалификационной работы - разработка и реализация мобильного приложения, предназначенного для совместного планирования поездок и снижения издержек групповой координации.

Для достижения поставленной цели решаются следующие задачи:
1. Проанализировать предметную область и выявить ограничения фрагментированного подхода к планированию поездок.
2. Рассмотреть существующие решения и определить их функциональные границы в контексте совместного планирования.
3. Сформулировать функциональный разрыв и обосновать целевую нишу разрабатываемого приложения.
4. Спроектировать архитектуру системы, модель данных и ключевые пользовательские сценарии.
5. Обосновать выбор методов и средств реализации.
6. Подготовить основу для программной реализации и последующей проверки сквозных сценариев.

Практическая значимость работы состоит в создании прикладного Android-приложения, в котором объединены пользовательски значимые функции: управление поездкой, совместный доступ, работа с идеями и маршрутом, учет расходов и балансов, погодный контекст, AI-предложения, уведомления и поддержка офлайн-работы с последующей синхронизацией [7][9].

Выпускная квалификационная работа включает три главы. В первой главе рассматриваются предметная область и существующие решения, во второй - проектирование системы, в третьей - программная реализация и результаты проверки сквозных пользовательских сценариев.
